% Generated draft: requires mathematical and source review.
\documentclass[11pt]{article}
\usepackage[margin=25mm]{geometry}
\usepackage{fontspec}
\setmainfont{Latin Modern Roman}
\usepackage{amsmath,amssymb,mathtools,unicode-math}
\setmathfont{Latin Modern Math}
\usepackage{enumitem,needspace,xurl,hyperref}
\setlength{\parindent}{0pt}
\setlength{\parskip}{5pt}
\setlength{\emergencystretch}{3em}
\newcommand{\N}{\mathbb{N}}
\newcommand{\Z}{\mathbb{Z}}
\newcommand{\Q}{\mathbb{Q}}
\newcommand{\R}{\mathbb{R}}
\newcommand{\C}{\mathbb{C}}
\newcommand{\F}{\mathbb{F}}
\newcommand{\A}{\mathbb{A}}
\newcommand{\PP}{\mathbb{P}}
\newcommand{\E}{\mathbb{E}}
\newcommand{\eps}{\varepsilon}
\newcommand{\dd}{\,\mathrm{d}}
\begin{document}
\section*{UnsolvedMath ID: 6700048}
\textbf{Source problem code:} catalogue entry
\subsection{Definitions and mathematical statement}
Let $(X^n,g)$ be a Riemannian manifold with
\[
 \operatorname{Sc}_g\ge\sigma>0.
\]
Fix positive parameters $\sigma_-,r,r_+$ satisfying $\sigma_-\le\sigma$ and $r_+\ge r$. The source asks when there exist a Riemannian $n$-manifold $(X^+,g^+)$ and an isometric embedding
\[
 \iota:X_{-r}\longrightarrow (X^+)_{-r_+}
\]
with $\operatorname{Sc}_{g^+}\ge\sigma_-$.

Here $X_{-r}$ and $(X^+)_{-r_+}$ denote the source's metric interiors. A usual boundary-distance convention would set
\[
 X_{-r}=\{x\in X:\operatorname{dist}_g(x,\partial X)>r\}.
\]
For an incomplete manifold with no specified manifold boundary, one instead needs a convention using the boundary in its metric completion. The record does not define which boundary, whether the inequality is strict, or the regularity of these interior sets. The same choice must be made for $X^+$.

Under the Riemannian interpretation, isometric embedding requires $\iota^*g^+=g$ on the indicated interior. If preservation of the ambient metric-space distance is intended, that is an additional convention requiring confirmation. The problem asks for conditions allowing the unchanged $r$-interior to lie at least the requested depth $r_+$ inside a new manifold with the lower scalar bound. It does not explicitly require $X^+$ to be complete, compact, to contain all of $X$, or to realize an extension of a prescribed topology. No such conditions are silently added, and no universal existence theorem is asserted.
\subsection{Short English statement}
When can a positively scalar-curved manifold's interior be placed isometrically deeper inside another manifold of the same dimension while allowing only the stated weakening of the scalar-curvature bound? The boundary/completion and embedding conventions must be fixed.

The supplied summary reports low-dimensional boundary-smoothing or extension constructions but describes the general arbitrary-dimensional problem as open. It supplies no exact theorem, hypotheses, or later reference. The OPEN-TRIAGE label and 2017 passage are retained with current status marked for review.
\subsection{Sources}
\begin{itemize}
\item Local UnsolvedMath registry, permanent ID 6700048, source code catalogue entry, scalar-curvature Question [?51], ``Extension Problem,'' proposed by Misha Gromov, 2017. Dataset curation and metadata: Ulam AI and UnsolvedMath Contributors, CC BY 4.0.
\item Misha Gromov, \emph{101 Questions, Problems and Conjectures around Scalar Curvature} (2017), printed tag [?51], occurrence 50, PDF page 57, \url{https://www.ihes.fr/~gromov/wp-content/uploads/2018/08/101-problemsOct1-2017.pdf}; normalized pdf-text extraction, accessed 2026-07-29.
\item No exact later extension theorem is supplied. Source-specific provenance: the author-hosted PDF is incomplete and publicly available; current status and statement-reuse terms NEEDS\_REVIEW, separately from the dataset licence.
\end{itemize}
\end{document}
